\section{Transformers}

\begin{frame}
    \frametitle{Previous Lecture}
    
    We already learn about following topics:
    \begin{itemize}
        \item Lab 9: Residual Networks (ResNets)
        \item Lab 7: Normalization (BatchNorm, LayerNorm, etc.)
        \item Lab 6: Deep Learning from Scratch
        \item Lab 0: Review Python and Jupyter Notebooks
    \end{itemize}

    Today we will implement the Transformers architecture. Transformers are just a top of the previous topics with some modifications!
\end{frame}

\begin{frame}{The one idea to hold onto}
  \begin{center}
    \alert{A GPT is just}\\[2mm]
    \large $\text{embedding} \;\to\; [\,\text{attention} + \text{MLP}\,]\times N \;\to\; \text{head} \;\to\; \text{softmax}$
  \end{center}
  \vfill
  \begin{columns}[T]
    \column{0.5\textwidth}
      \textbf{You already built this}
      \begin{itemize}
        \item tensors / vectors
        \item \texttt{loss.backward()}, \texttt{optimizer.step()}
        \item the training loop
        \item \texttt{softmax} + cross-entropy
      \end{itemize}
    \column{0.5\textwidth}
      \textbf{Genuinely new today}
      \begin{itemize}
        \item[\textcolor{unistemerald}{$\bigstar$}] \alert{attention}
      \end{itemize}
      \vspace{4mm}
      {\footnotesize Everything else is something from your last labs.}
  \end{columns}
\end{frame}

\begin{frame}{What you did last time $\to$ what it becomes today}
  \begin{center}
  \renewcommand{\arraystretch}{1.4}
  \begin{tabular}{l@{\quad$\Rightarrow$\quad}l}
    \textbf{Already in your toolbox} & \textbf{Today} \\\hline
    \texttt{softmax} + cross-entropy & predict the \emph{next token} = classification \\
    \texttt{loss.backward()} (a black box) & you'll \emph{read} the 14 lines that are it \\
    a few \texttt{nn.Linear} + ReLU & the MLP inside every block \\
    the training loop & \emph{unchanged} — it trains a GPT \\
  \end{tabular}
  \end{center}
  \vfill
  \begin{center}\footnotesize So: only \alert{attention} is new. Let's earn it.\end{center}
\end{frame}

% --- Embedding ---
\begin{frame}
    \frametitle{Embedding}

    \begin{definition}[Embedding]
        An \textbf{embedding} is a mapping from objects --- from words, molecules, even real world objects --- to vectors that we can do calculations with. The embedding is a learned function, and the vectors it produces are called \textbf{embeddings} of the original objects.
    \end{definition}

    \pause

    \textcolor{red}{\textbf{Question:}} What is the \textbf{good} embedding?

    \pause

    To be honest, we don't know. But we can say that a good embedding is one that \textbf{makes the task easier} for the model. For example, if we are doing language modeling, a good embedding would be one that captures the meaning of words in a way that makes it easier for the model to predict the next word.
\end{frame}

\begin{frame}
    \frametitle{Embedding: a Picture}

    \begin{columns}[T,onlytextwidth]
        \column{0.42\textwidth}
            Each word becomes a \alert{point in space}.
            \vspace{2mm}
            \begin{itemize}
                \item \texttt{"cat"} $\;\mapsto\; (1.6,\, 1.8)$
                \item \texttt{"car"} $\;\mapsto\; (1.9,\, 0.4)$
            \end{itemize}
            \vspace{2mm}
            Words with similar \emph{meaning} end up \alert{close together}.
            \vspace{2mm}

            {\footnotesize (Here $d=2$ so we can draw it. In a real GPT, $d$ is hundreds or thousands.)}
            \vspace{2mm}

            \pause

            \textcolor{red}{\textbf{Question:}} where would \texttt{"puppy"} go? What about \texttt{"taxi"}?
        \column{0.55\textwidth}
            \begin{center}
            \begin{tikzpicture}[scale=1.9, font=\scriptsize, >=Stealth]
                % axes
                \draw[->, gray] (0,0) -- (2.6,0) node[below] {$x_1$};
                \draw[->, gray] (0,0) -- (0,2.4) node[left] {$x_2$};
                % animal cluster
                \fill[grn] (1.6,1.8) circle (1.4pt) node[above] {cat};
                \fill[grn] (1.3,2.0) circle (1.4pt) node[above] {kitten};
                \fill[grn] (1.9,1.6) circle (1.4pt) node[right] {dog};
                \draw[grn, dashed] (1.6,1.8) ellipse (0.55 and 0.42);
                % vehicle cluster
                \fill[oge] (1.9,0.4) circle (1.4pt) node[below] {car};
                \fill[oge] (2.3,0.6) circle (1.4pt) node[right] {truck};
                \fill[oge] (1.7,0.7) circle (1.4pt) node[above left] {bus};
                \draw[oge, dashed] (2.0,0.55) ellipse (0.55 and 0.42);
                % token -> vector arrow
                \node[draw=unistnavy, fill=gra, rounded corners=1pt, inner sep=2pt] (tok) at (0.5,1.6) {\texttt{"cat"}};
                \draw[->, unistnavy, thick] (tok) -- (1.55,1.78);
            \end{tikzpicture}
            \end{center}
    \end{columns}
\end{frame}

% Similarity

\begin{frame}
    \frametitle{Similarity in Embedding Space}

    Now words are vectors. \textcolor{red}{\textbf{Question:}} how do we measure that \texttt{"cat"} is closer to \texttt{"kitten"} than to \texttt{"car"}?

    \pause

    \begin{definition}[Similarity]
        \textbf{Similarity} in embedding space is a measure of how close two embeddings are. We compute it with the \alert{inner product}:
        \[
            \mathrm{sim}(u, v) = \langle u, v \rangle = u^\top v = \sum_{k=1}^{d} u_k v_k .
        \]
    \end{definition}

    \pause

    One number per pair: \alert{large} $\Rightarrow$ similar, \alert{small (or negative)} $\Rightarrow$ not similar.\footnote{You may have heard of \emph{cosine similarity} --- that is just the inner product after normalizing the lengths: $\langle u, v\rangle = \|u\|\|v\|\cos\theta$. We stick to the plain inner product today.}
\end{frame}

\begin{frame}
    \frametitle{Similarity: a Picture}

    \begin{columns}[T,onlytextwidth]
        \column{0.42\textwidth}
            Same picture as before, but now look at the \alert{angles}.
            \vspace{2mm}
            \begin{itemize}
                \item \textcolor{grn}{cat $\cdot$ kitten}: small angle $\Rightarrow$ \alert{large} inner product
                \item \textcolor{oge}{cat $\cdot$ car}: large angle $\Rightarrow$ \alert{small} inner product
            \end{itemize}
            {\footnotesize (Here all vectors have about the same length, so the angle is what's left to vary.)}
            \vspace{2mm}

            \pause

            \textcolor{red}{\textbf{Question:}} we have $n$ words. How many pairs do we need to compare?

            \pause

            All $n \times n$ of them... do we write a double \texttt{for} loop?
        \column{0.55\textwidth}
            \begin{center}
            \begin{tikzpicture}[scale=1.9, font=\scriptsize, >=Stealth]
                % axes
                \draw[->, gray] (0,0) -- (2.6,0) node[below] {$x_1$};
                \draw[->, gray] (0,0) -- (0,2.4) node[left] {$x_2$};
                % vectors
                \draw[->, grn, thick] (0,0) -- (1.6,1.8) node[above] {cat};
                \draw[->, grn, thick] (0,0) -- (1.3,2.0) node[above] {kitten};
                \draw[->, oge, thick] (0,0) -- (1.9,0.4) node[right] {car};
                % angle marks
                \draw[grn] (0.45,0.51) arc[start angle=48.4, end angle=57, radius=0.68];
                \node[grn] at (0.62,0.78) {\tiny $\theta_{\text{small}}$};
                \draw[oge] (0.50,0.105) arc[start angle=11.9, end angle=48.4, radius=0.51];
                \node[oge] at (0.72,0.30) {\tiny $\theta_{\text{large}}$};
            \end{tikzpicture}
            \end{center}
    \end{columns}
\end{frame}


\begin{frame}
    \frametitle{All Pairs at Once: Matrix Multiplication}

    No \texttt{for} loop! Stack the $n$ embeddings as \alert{rows} of a matrix $X \in \mathbb{R}^{n \times d}$. Then
    \[
        S = X X^\top \in \mathbb{R}^{n \times n},
        \qquad
        S_{ij} = x_i^\top x_j = \mathrm{sim}(x_i, x_j).
    \]

    \pause

    \begin{center}
    \begin{tikzpicture}[scale=1.5, font=\scriptsize]
        % X (n x d)
        \draw[unistnavy, thick] (0,0) rectangle (0.9,1.2);
        \fill[grn!40] (0,0.6) rectangle (0.9,0.9);
        \node at (0.45,0.75) {\tiny $x_i^\top$};
        \node at (0.45,1.4) {$X$};
        \node at (0.45,-0.2) {\tiny $n \times d$};
        % times
        \node at (1.25,0.6) {$\times$};
        % X^T (d x n)
        \draw[unistnavy, thick] (1.6,0.15) rectangle (2.8,1.05);
        \fill[oge!40] (2.2,0.15) rectangle (2.5,1.05);
        \node at (2.35,0.6) {\tiny $x_j$};
        \node at (2.2,1.4) {$X^\top$};
        \node at (2.2,-0.2) {\tiny $d \times n$};
        % equals
        \node at (3.15,0.6) {$=$};
        % S (n x n)
        \draw[unistnavy, thick] (3.5,0) rectangle (4.7,1.2);
        \fill[grn!25] (3.5,0.6) rectangle (4.7,0.9);
        \fill[oge!25] (4.1,0) rectangle (4.4,1.2);
        \fill[unistemerald!70] (4.1,0.6) rectangle (4.4,0.9);
        \node at (4.25,0.75) {\tiny $S_{ij}$};
        \node at (4.1,1.4) {$S$};
        \node at (4.1,-0.2) {\tiny $n \times n$};
    \end{tikzpicture}
    \end{center}

    \pause

    \textcolor{red}{\textbf{Question:}} what does row $i$ of $S$ tell us?

    \pause

    How similar word $i$ is to \alert{every} word --- including itself. Remember this row; it is about to become \textbf{attention}.
\end{frame}

% --- Attention ---

\begin{frame}{Attention in Language}
  \begin{quote}
    ``The cat that the dog chased \underline{\hspace{1.2cm}}\,''
  \end{quote}
  \begin{itemize}
    \item To fill the blank you must look \alert{back} at the \emph{right} earlier word.
    \item An MLP has a fixed input and \alert{no sense of order} or of which earlier word matters.
    \item We need a way to \alert{mix information across positions, based on content}.
    \item[] \hfill $\Rightarrow$ that mechanism is \textbf{attention}.
  \end{itemize}
\end{frame}

\begin{frame}{Attention in Vision}
    \begin{center}
        \includegraphics[width=0.8\textwidth]{src/overthinking.jpg}
    \end{center}\footnote{Image from \url{https://www.boredpanda.com/captcha-struggles-fails/}}
\end{frame}

\begin{frame}
    \frametitle{From Similarity to Attention}

    \textcolor{red}{\textbf{Question:}} we already have a similarity matrix $S = XX^\top$. Why not use it to decide \emph{which earlier words matter}?

    \pause

    That is almost exactly attention --- with one twist. Each word plays \alert{three roles}, each a learned linear map of $x$:
    \begin{center}
    \renewcommand{\arraystretch}{1.3}
    \begin{tabular}{lll}
        \textbf{Query} & $Q = X W_Q$ & ``what am I \emph{looking for}?'' \\
        \textbf{Key}   & $K = X W_K$ & ``what do I \emph{offer}?'' \\
        \textbf{Value} & $V = X W_V$ & ``what do I \emph{carry}?'' \\
    \end{tabular}
    \end{center}

    \pause

    Same picture as before, just with $Q$ and $K$ instead of $X$ and $X^\top$:
    \[
        S = Q K^\top \in \mathbb{R}^{n \times n}
        \qquad \text{(still just inner products of all pairs!)}
    \]
\end{frame}

\begin{frame}
    \frametitle{Remark: Attention Measures Similarity in a New Basis}

    \textcolor{red}{\textbf{Question:}} what does multiplying by $W_Q$ actually \emph{do} to an embedding?

    \pause

    \begin{columns}[T,onlytextwidth]
        \column{0.55\textwidth}
            A linear map \alert{recombines the basis}: each new axis is a mix of the old ones. So
            \[
                Q K^\top = (X W_Q)(X W_K)^\top = X \underbrace{\left(W_Q W_K^\top\right)}_{\text{learned!}} X^\top
            \]
            Attention does \emph{not} compare embeddings in their original coordinates --- it compares them in a \alert{learned, tilted basis} where the axes of the two words \alert{cross}.

            \vspace{2mm}
            \pause

            The model gets to \emph{choose which directions of meaning matter} for matching. With $W_Q = W_K = I$ we would be back to plain $XX^\top$.
        \column{0.42\textwidth}
            \begin{center}
            \begin{tikzpicture}[scale=1.3, font=\scriptsize, >=Stealth]
                % original basis (gray)
                \draw[->, gray, thick] (0,0) -- (1.6,0) node[below] {$e_1$};
                \draw[->, gray, thick] (0,0) -- (0,1.6) node[left] {$e_2$};
                % learned basis (tilted, emerald)
                \draw[->, unistemerald, very thick] (0,0) -- (1.45,0.65) node[right] {$W e_1$};
                \draw[->, unistemerald, very thick] (0,0) -- (-0.55,1.5) node[above] {$W e_2$};
                % same point, two coordinate readings
                \fill[oge] (0.9,1.1) circle (1.4pt) node[above right] {cat};
                \draw[gray, dashed] (0.9,1.1) -- (0.9,0);
                \draw[gray, dashed] (0.9,1.1) -- (0,1.1);
            \end{tikzpicture}\\[1mm]
            {\scriptsize same embedding, read off \\ in two \alert{crossing} bases}
            \end{center}
    \end{columns}
\end{frame}

\begin{frame}
    \frametitle{Attention: the Full Formula}

    \begin{definition}[Attention]
        \[
            \mathrm{Attention}(Q, K, V) = \underbrace{\mathrm{softmax}\!\left(\frac{Q K^\top}{\sqrt{d}}\right)}_{\text{weights, each row sums to } 1} V
        \]
    \end{definition}

    \pause

    Read it from the inside out:
    \begin{enumerate}
        \item $QK^\top$: similarity of every pair --- you built this already.
        \item $\mathrm{softmax}$ (row-wise): turn similarities into weights --- you built this already too!
        \item $\cdots \times V$: each word's output is a \alert{weighted average of the values}.
    \end{enumerate}

    \pause

    \textcolor{red}{\textbf{Question:}} so what is genuinely new here?

    \pause

    Only the idea of $Q$, $K$, $V$. The rest is matrix multiplication and softmax.
\end{frame}

\begin{frame}
    \frametitle{Three Small Details (Brief!)}

    \begin{columns}[T,onlytextwidth]
        \column{0.58\textwidth}
            \textbf{1. Scale by $\sqrt{d}$.} When $d$ is large, inner products get large, and softmax \alert{saturates}. Dividing by $\sqrt{d}$ keeps it in a healthy range. That's it.

            \vspace{2mm}
            \pause

            \textbf{2. Causal masking.} GPT predicts the \emph{next} word, so position $i$ must not peek at the \alert{future} ($j > i$). We set those scores to $-\infty$ before softmax $\Rightarrow$ weight $0$.

            \vspace{2mm}
            \pause

            \textbf{3. Multi-head.} Instead of one attention with width $d$, run $h$ \alert{smaller ones in parallel} (each width $d/h$) and concatenate. Each head can look for a different relationship.
        \column{0.38\textwidth}
            \begin{center}
            \begin{tikzpicture}[scale=0.55, font=\tiny]
                % causal mask grid 4x4 (row 0 at top)
                \foreach \i in {0,...,3} {
                    \foreach \j in {0,...,3} {
                        \pgfmathparse{\j > \i ? 1 : 0}
                        \ifnum\pgfmathresult>0
                            \fill[gray!30] (\j,-\i) rectangle ++(1,-1);
                            \node at (\j+0.5,-\i-0.5) {$-\infty$};
                        \else
                            \fill[unistemerald!30] (\j,-\i) rectangle ++(1,-1);
                            \node at (\j+0.5,-\i-0.5) {$S_{\i\j}$};
                        \fi
                    }
                }
                \draw[unistnavy] (0,0) grid[step=1] (4,-4);
                \draw[unistnavy, thick] (0,0) rectangle (4,-4);
                \node[anchor=south] at (2,0.1) {keys $j \rightarrow$};
                \node[anchor=south, rotate=90] at (-0.3,-2) {queries $i \rightarrow$};
            \end{tikzpicture}\\[1mm]
            {\scriptsize the causal mask: \\ gray = ``no peeking''}
            \end{center}
    \end{columns}
\end{frame}

\begin{frame}
    \frametitle{See It on a Real Model (Live)}

    \begin{columns}[c]
        \column{0.33\textwidth}
            \begin{center}
                \qrcode[height=2.2cm]{https://poloclub.github.io/transformer-explainer/}\\[1mm]
                {\footnotesize \textbf{Transformer Explainer}\\ real GPT-2 in your browser}
            \end{center}
        \column{0.33\textwidth}
            \begin{center}
                \qrcode[height=2.2cm]{https://bbycroft.net/llm}\\[1mm]
                {\footnotesize \textbf{bbycroft.net/llm}\\ 3D data-flow walkthrough}
            \end{center}
        \column{0.33\textwidth}
            \begin{center}
                \qrcode[height=2.2cm]{https://www.3blue1brown.com/lessons/attention}\\[1mm]
                {\footnotesize \textbf{3Blue1Brown}\\ attention, animated}
            \end{center}
    \end{columns}

    \vfill
    \begin{center}\footnotesize Type a sentence, change a word, watch the attention weights move.\end{center}
\end{frame}